\documentclass[conference]{IEEEtran}
\IEEEoverridecommandlockouts
% The preceding line is only needed to identify funding in the first footnote. If that is unneeded, please comment it out.
\usepackage{cite}
\usepackage{amsmath,amssymb,amsfonts}
\usepackage{algorithmic}
\usepackage{graphicx}
\usepackage{textcomp}
\usepackage{xcolor}
\usepackage{float}
\usepackage{subcaption}
\usepackage{algorithm}
\usepackage{algpseudocode}


% ====== 新增以下這段來支援中文 ======
\usepackage[AutoFakeBold=true, AutoFakeSlant=true]{xeCJK}\setmainfont{Times New Roman}
% 設定主要的繁體中文字體，您可以根據您的作業系統或編譯環境更改
%\setCJKmainfont{Noto Sans TC} % Windows 常用的字體
\setCJKmainfont{標楷體}   % 如果想要比較正式的襯線字體

\def\BibTeX{{\rm B\kern-.05em{\sc i\kern-.025em b}\kern-.08em
    T\kern-.1667em\lower.7ex\hbox{E}\kern-.125emX}}
\begin{document}

\title{Memetic Algorithm with NEH Seeding and Tabu-based Local Search for Permutation Flowshop Scheduling\\
{\large 結合 NEH 初始化與 Tabu 區域搜尋之迷因演算法於排列流水線排程問題之研究}
}


\author{\IEEEauthorblockN{1\textsuperscript{st} Hung Ming-Chun}
\IEEEauthorblockA{\textit{Department of Computer Science and} \\
\textit{Information Engineering}\\
\textit{National Taiwan Normal University}\\
Taipei, Taiwan(R.O.C) \\
61447026s@ntnu.edu.tw}
\and
\IEEEauthorblockN{2\textsuperscript{nd} Hsieh Yao-yu}
\IEEEauthorblockA{\textit{Graduate Institute of} \\
\textit{Science Education}\\
\textit{National Taiwan Normal University}\\
Taipei, Taiwan(R.O.C) \\
61345001s@ntnu.edu.tw}
\and
\IEEEauthorblockN{3\textsuperscript{rd} Given Name Surname}
\IEEEauthorblockA{\textit{Department of Computer Science and} \\
\textit{Information Engineering}\\
\textit{National Taiwan Normal University}\\
Taipei, Taiwan(R.O.C) \\
@ntnu.edu.tw}
}



\maketitle

\begin{abstract}

\end{abstract}

\section{Introduction}

排列流水線排程問題（Permutation Flowshop Scheduling Problem, PFSP）是排程領域中經典且具高度實務價值的組合最佳化問題，廣泛應用於製造系統、生產排程與資源配置等實務問題中。此問題的目標是在多台機器上決定所有工作的加工順序，且所有機器需依相同順序處理工作，使總完工時間（makespan）最小化。由於 PFSP 在一般情況下屬於 NP-hard 問題，當工作數與機器數增加時，傳統精確演算法往往難以在合理時間內求得最佳解。\\

基因演算法（Genetic Algorithm, GA）因具備良好的全域搜尋能力與彈性的排列編碼方式，常被用於求解 PFSP。然而，傳統 GA 在演化後期容易因族群多樣性下降而產生早熟收斂，使得搜尋停滯於區域最佳解附近。為了同時兼顧全域探索能力與局部開發能力，本研究採用迷因演算法（Memetic Algorithm, MA）的觀點，以 GA 作為族群式搜尋主體，並結合軌跡式元啟發法作為強化局部搜尋能力的改良機制。\\

本研究使用 permutation-based encoding 表示解，並設計一個特殊初始化程序與軌跡式元啟發法的整合架構。具體而言，我們首先以 NEH 建構式啟發法產生高品質初始種子，以改善初始族群品質；接著比較多種適用於排列問題的交配算子與突變設定；最後導入以 swap 為鄰域的 Tabu-based local search，針對菁英個體進行局部強化，形成一個用於 PFSP 的 memetic framework。\\

本研究的目的不僅在於建構一個有效的 PFSP 求解方法，同時亦系統性分析不同設計選擇對演算法效能之影響，包括交配算子、突變率、突變方式、NEH 初始化，以及 Tabu-based local search 的整合方式。實驗將使用 Taillard benchmark 測資，並透過多次獨立執行比較各方法在解品質、穩定性與運算成本上的差異。\\



%%%%%%%%%%%%%%%%%%%%%%%%%%%%%%%%%%%%%%%%%%%%%%%%%%

\section{Problem definition}
%這邊著重在介紹演算法，實作細節在Implement Detail
本研究考慮排列流水線排程問題（Permutation Flowshop Scheduling Problem, PFSP）。設共有 \(n\) 個工作與 \(m\) 台機器，所有工作皆須依固定順序經過各台機器加工，即每個工作都必須依序通過 \(M_1, M_2, \dots, M_m\)。此外，每台機器在任一時間僅能處理一個工作，且加工過程不可中斷。

PFSP 的特徵在於所有機器皆以相同的工作順序進行加工。換言之，若某一解表示為一個工作排列 \(\pi = (\pi_1, \pi_2, \dots, \pi_n)\)，則此排列同時決定所有機器上的加工順序。因此，本問題的核心在於尋找一組最佳排列，使所有工作完成加工的最終時間最小。

令 \(p_{i,\pi_k}\) 表示工作 \(\pi_k\) 在機器 \(M_i\) 上的加工時間，\(C(\pi_k, i)\) 表示工作 \(\pi_k\) 在機器 \(M_i\) 上的完工時間，則其遞迴關係可表示為
\[C(\pi_k, i) = \max \{ C(\pi_{k-1}, i),\ C(\pi_k, i-1) \} + p_{i,\pi_k}.\]
其中，\(C(\pi_0, i) = 0\) 與 \(C(\pi_k, 0) = 0\)。本研究之目標函數為最小化總完工時間（makespan）：
\[C_{\max} = C(\pi_n, m).\]

因此，PFSP 可視為一個在所有可行工作排列中，尋找使 \(C_{\max}\) 最小之最佳排列的組合最佳化問題。

%%%%%%%%%%%%%%%%%%%%%%%%%%%%%%%%%%%%%%%%%%%%%%%%%%
\section{Proposed Memetic Algorithm}
\subsection{Solution Representation}
本研究採用排列式編碼表示解。若共有 \(n\) 個工作，則一個候選解可表示為一個長度為 \(n\) 的排列向量，其中每個元素對應一個唯一的工作編號，且每個工作僅出現一次。

例如，當一個解表示為 \((3,6,4,1,5,2)\) 時，代表所有機器皆依照工作 \(3 \rightarrow 6 \rightarrow 4 \rightarrow 1 \rightarrow 5 \rightarrow 2\) 的順序進行加工。由於 PFSP 要求所有機器採用相同的工作順序，因此此種排列式編碼方式可直接對應到一個合法排程，並適合搭配排列問題常見的交配、突變與區域搜尋運算子。
\subsection{Fitness Evaluation}
本研究以總完工時間作為解的適應值評估指標。對於任一給定的工作排列，系統會依序模擬各工作在每台機器上的加工流程，並計算最後一個工作於最後一台機器上的完工時間，作為該解的 makespan。

在實作上，我們設計了一個排程評估函數。令 \texttt{machine\_end\_time[i]} 表示機器 \(M_i\) 的最早可用時間，當某一工作進入機器 \(M_i\) 時，其開始時間取決於該工作在前一台機器上的完工時間與該機器目前最早可用時間兩者中的較大值。完成所有工作與所有機器的模擬後，最後一台機器的完工時間即為該排列的 makespan。
\subsection{Genetic Framework}

本研究以基因演算法作為迷因演算法的主體架構。演算法維護一組族群，並透過選擇、交配與突變逐代演化，以持續改善解的品質。

在親代選擇方面，本研究採用錦標賽選擇法。每次由族群中隨機挑選若干個體，並選出其中 makespan 最小者作為親代，以兼顧選擇壓力與族群多樣性。在世代更新方面，本研究採用菁英保留策略，將當代最佳個體直接保留至下一代，以避免優良解在交配或突變過程中流失。

此外，本研究以目標函數評估次數（function evaluations, FFE）作為演算法的主要搜尋預算控制方式。每當候選解被評估一次時，FFE 計數器即增加 1，並於達到預設上限時終止演化流程。
\subsection{Crossover and Mutation Operators}
由於本研究採用排列式編碼表示解，傳統的單點或多點交配無法直接套用，否則容易產生重複工作或遺漏工作的非法排列。因此，本研究實作了四種常見的排列交配算子，分別為順序交配（Order Crossover, OX）、線性順序交配（Linear Order Crossover, LOX）、部分映射交配（Partially Mapped Crossover, PMX）以及循環交配（Cycle Crossover, CX），並比較其在 PFSP 上的效能表現。

OX 會先保留父代的一段連續基因片段，再依照另一個父代中的相對出現順序填補其餘空位。此方法能在保留部分局部順序資訊的同時，維持子代為合法排列。

LOX 與 OX 相似，同樣先保留一段父代片段，但其餘位置的填補方式採用由左至右的線性順序，而非循環繞回。此特性使其更強調線性結構的保存。

PMX 則透過兩個切點建立父代片段間的對應關係，並利用映射規則修正其餘位置的衝突，以同時兼顧絕對位置與相對結構的保留。

CX 的核心概念是在父代之間找出循環結構，使子代各位置上的基因來自其中一個父代在相同位置上的基因，因此特別強調絕對位置資訊的保留。

在突變方面，本研究考慮兩種排列問題常見的突變方式：交換突變(Swap)與插入突變(Insertion)。交換突變會隨機選擇兩個位置並交換其工作；插入突變則會抽出一個工作並插入到另一個位置。透過比較不同突變方式與突變率設定，可分析其對族群多樣性與局部搜尋能力的影響。

\subsection{NEH-based Initialization}
本研究在初始族群中引入 NEH 建構式啟發法所產生的高品質種子解。NEH 方法首先計算各工作在所有機器上的總加工時間，並依總加工時間由大至小排序；接著依序將工作插入目前部分序列的最佳位置，以逐步建構一組品質較佳的初始解。

在本研究中，初始族群並非完全隨機生成，而是於族群中注入一個由 NEH 產生的種子個體，其餘個體則維持隨機排列。此設計可在保留族群多樣性的同時，提升演化初期的搜尋品質與收斂效率。

藉由這種逐一探索局部最佳解的插入方式，NEH 能夠產生結構極佳的初始解。將此 NEH 解注入基因演算法的初始族群中，能確保演算法從一個具備優良基因片段的起點開始演化，進而顯著提升最終搜尋到全局最佳解的機率與效率。

\subsection{Tabu-based Local Search}
為了提升演算法的局部開發能力，本研究於基因演化流程中加入 Tabu-based 區域搜尋，形成迷因演算法架構。區域搜尋以交換鄰域為基礎，亦即對目前解中任兩個位置的工作進行交換，以產生鄰近解。

在每一世代完成族群更新後，本研究選取當代最佳個體作為區域搜尋的對象，並對其執行固定步數的禁忌搜尋。搜尋過程中，系統以交換的工作對作為禁忌移動記錄，避免搜尋在短時間內重複走回已探索區域；若某禁忌移動可產生優於目前歷史最佳解的結果，則透過特赦準則予以接受。區域搜尋完成後，改良後的個體會回寫至族群中，並作為後續世代演化的基礎。

此設計使演算法能以 GA 維持全域探索能力，並以禁忌搜尋加強對高品質個體的局部精修，藉此兼顧 exploration 與 exploitation。

綜合而言，本研究所提出的迷因演算法以基因演算法作為全域搜尋主體，並結合 NEH 初始化與 Tabu-based 區域搜尋，以同時提升演化初期的解品質與後期的局部開發能力。此架構兼顧探索與開發，作為後續實驗比較的主要方法。

\subsection{Population-level Tabu Mechanism}

除了將 Tabu Search 用於菁英個體的局部改良外，本研究亦曾於族群演化過程中加入一個 population-level tabu mechanism，用以避免短時間內重複產生相同或過於相似的排列。具體而言，本機制會將已產生之個體以雜湊值形式記錄，並於 offspring 生成時檢查其是否已存在於當前族群或近期 tabu list 中。若某候選個體被判定為 tabu，則該個體不會直接加入下一代，而是重新嘗試產生新的 offspring；若多次嘗試後仍無法得到合法且非 tabu 的候選解，則以隨機個體補入族群。

此機制的目的在於降低族群中重複解的比例，並提升搜尋過程中的多樣性。然而，與針對個體進行局部精修的 Tabu-based local search 不同，此處之 tabu 機制主要作用於族群層級，屬於一種 diversity control strategy，而非局部搜尋程序本身。
%%%%%%%%%%%%%%%%%%%%%%%%%%%%%%%%%%%%%%%%%%%%%%%%%%%
\newpage
\section{Experiment Design}
\fontsize{10pt}{12pt}\selectfont

\subsection{Experimental Settings}

本研究使用 Taillard benchmark 作為測試資料集，並對每一組演算法設定進行 20 次獨立執行，以降低隨機因素對結果的影響。除非另有說明，各組實驗皆採用相同的基本設定，以確保比較的公平性。

在演化參數方面，本研究設定族群大小為 50，交配率為 0.8，突變率則依實驗目的設定為不同數值。親代選擇採用錦標賽選擇法(k-tournament)，其錦標賽大小設定為 3，並於每一世代保留當代最佳個體。對於區域搜尋部分，Tabu-based local search 之最大搜尋步數與禁忌期限則依預設設定執行。

此外，本研究以目標函數評估次數(FFE)作為主要停止條件，以確保不同方法之比較具有公平性。設族群大小為 50，則本研究將最大 FFE 設定為 50000。此外，由於加入區域搜尋後，每一世代所需的評估數可能顯著增加，因此本研究亦額外記錄實際完成的世代數，作為分析演算法搜尋成本與演化效率的輔助指標。
\subsection{Comparison Plan}

為了系統性分析各設計選擇對演算法效能的影響，本研究依序進行以下比較。首先，比較不同交配算子在相同突變與初始化設定下的表現，以選出較適合作為後續實驗基礎的交配方式。其次，在固定交配算子的情況下，比較不同突變率與突變方式，以分析其對搜尋品質與穩定性的影響。

在決定較佳的基因演算法設定後，本研究進一步比較是否使用 NEH 初始化，以分析特殊初始化程序對演化初期品質與最終結果的影響。最後，在前述較佳設定下加入 Tabu-based local search，並與未使用區域搜尋的版本進行比較，以驗證迷因演算法架構是否能進一步提升解品質。

\subsection{Performance Measures}

本研究在比較不同方法時，記錄各組實驗於 20 次獨立執行下的最佳值、平均值、最差值與標準差。此外，亦統計平均評估次數與平均執行時間，以評估各方法之運算成本。若可取得最佳已知解（best known solution, BKS），亦進一步計算平均偏差百分比，以衡量解品質與最佳已知解之差距。
若已知某測資之BKS，則本研究進一步計算平均偏差百分比（average deviation percentage）：
\[\text{Avg. Dev. (\%)} = \frac{\text{Avg} - \text{BKS}}{\text{BKS}} \times 100\%.\]\par

\subsection{Gnuplot Visualizer}
為了直觀分析演算法的收斂行為與跳脫區域最佳解的過程，本實作內建了基於 \texttt{Gnuplot} 的視覺化工具。程式透過 \texttt{popen} 開啟 \texttt{gnuplot} 的行程通訊（Pipe），並利用標準輸入字元 \texttt{'-'} 將 C++ 陣列中的 \texttt{best\_makespan\_history} 與 \texttt{makespan\_history} 即時傳遞給繪圖引擎。

在大量平行測試的環境下，系統會透過 \texttt{save\_plot} 函數在背景靜默生成圖表。該函數會自動切換 Gnuplot 的 \texttt{terminal} 為 \texttt{jpeg}，並根據演算法名稱、測資名稱及最佳解數值自動命名檔案存入指定目錄，確保實驗數據的可重現性與易讀性。

\section{Results and Discussion}
\subsection{Effect of Crossover Operator}
為了分析不同交配算子對演算法效能之影響，本實驗固定族群大小、交配率、突變率與突變方式，並僅改變交配算子種類。具體而言，本實驗固定使用插入突變，突變率設為 0.1，且暫不啟用 NEH 初始化與 Tabu-based local search，以公平比較 OX、LOX、PMX 與 CX 四種交配算子之表現。

表~\ref{tab:crossover_avg} 彙整了不同交配算子在九組 Taillard 測資上的平均 makespan。若以平均解品質作為主要比較指標，LOX 在多數測資中取得較佳結果，整體表現最為穩定。PMX 與 OX 在部分測資中亦能取得接近 LOX 的結果，而 CX 雖然在若干案例中也具有競爭力，但整體平均表現仍略遜於 LOX。

綜合而言，LOX 在解品質與穩定性之間取得較佳平衡，因此本研究後續實驗將以 LOX 作為主要交配算子，並在此基礎上進一步比較突變率、突變方式、NEH 初始化以及 Tabu-based local search 的效果。包含最小值、最大值與標準差的完整結果整理於附錄表~\ref{tab:crossover_full}。

\begin{table}[H]
\centering
\caption{在不同交配算子下的平均 makespan}
\label{tab:crossover_avg}
\begin{tabular}{c|c|c|c|c}
\hline
Instance & OX & LOX & PMX & CX \\
\hline
tai20\_5\_1   & 1295.85 & 1296.05 & \textbf{1294.60} & 1297.10 \\
tai20\_10\_1  & 1627.55 & \textbf{1621.85} & 1623.40 & 1629.55 \\
tai20\_20\_1  & 2373.00 & 2376.05 & \textbf{2360.80} & 2363.25 \\
tai50\_5\_1   & 2742.85 & \textbf{2741.35} & 2742.80 & 2744.75 \\
tai50\_10\_1  & \textbf{3145.70} & 3152.65 & 3150.80 & 3150.05 \\
tai50\_20\_1  & 4056.80 & \textbf{4052.85} & 4054.40 & 4057.50 \\
tai100\_5\_1  & \textbf{5512.90} & 5517.20 & 5516.15 & 5513.15 \\
tai100\_10\_1 & 5963.00 & \textbf{5939.00} & 5970.60 & 5959.85 \\
tai100\_20\_1 & 6687.35 & 6685.45 & \textbf{6674.90} & 6679.70 \\
\hline
\end{tabular}
\end{table}

\subsection{Effect of Mutation Settings}
在決定以 LOX 作為主要交配算子後，本研究進一步分析突變設定對演算法效能之影響。考量突變率與突變方式之間可能存在交互作用，本研究同時比較插入突變與交換突變兩種方式，並搭配兩種突變率（0.1 與 0.5）進行實驗。其餘參數皆維持一致，且暫不啟用 NEH 初始化與 Tabu-based local search。

表~\ref{tab:mutation_settings} 彙整了四種突變設定在九組 Taillard 測資上的平均 makespan。由結果可觀察到，突變率對整體表現具有顯著影響。在多數測試案例中，較高的突變率（0.5）在插入突變下能取得最佳平均解品質，顯示較強的擾動有助於維持族群多樣性並避免搜尋過早收斂。

綜合而言，插入突變搭配突變率 0.5 在本研究中展現最佳整體表現，因此後續實驗將採用此設定作為預設突變策略。
\begin{table}[H]
\centering
\caption{在不同突變組合設定下的平均 makespan}
\label{tab:mutation_settings}
\begin{tabular}{c|c|c|c|c}
\hline
Instance & Insertion + 0.1 & Insertion + 0.5 & Swap + 0.1 & Swap + 0.5 \\
\hline
tai20\_5\_1   & 1293.35 & \textbf{1291.30} & 1297.80 & 1297.00 \\
tai20\_10\_1  & 1624.35 & \textbf{1603.65} & 1640.80 & 1617.65 \\
tai20\_20\_1  & 2358.40 & \textbf{2341.05} & 2381.40 & 2346.30 \\
tai50\_5\_1   & 2743.70 & 2734.80 & 2743.75 & \textbf{2729.20} \\
tai50\_10\_1  & 3145.95 & \textbf{3116.25} & 3157.45 & 3139.20 \\
tai50\_20\_1  & 4073.40 & \textbf{4011.25} & 4089.00 & 4051.35 \\
tai100\_5\_1  & 5511.75 & \textbf{5496.10} & 5512.20 & 5506.80 \\
tai100\_10\_1 & 5962.25 & \textbf{5859.90} & 5993.90 & 5906.55 \\
tai100\_20\_1 & 6685.95 & \textbf{6579.90} & 6690.70 & 6581.90 \\
\hline
\end{tabular}
\end{table}

\subsection{Effect of NEH Initialization}
在確定交配算子與突變設定後，本研究進一步分析 NEH 初始化對演算法效能之影響。為了公平比較，本實驗固定使用 LOX 作為交配算子，並採用插入突變搭配突變率 0.5；此外，Tabu-based local search 暫不啟用。在此設定下，比較是否使用 NEH 初始化對解品質的影響。

表~\ref{tab:neh} 彙整了在有無 NEH 初始化下，各測試案例的平均 makespan。由結果可觀察到，在大多數測試案例中，使用 NEH 初始化皆能有效提升解品質，且在中大型問題（如 50-job 與 100-job）上改善特別明顯。例如在 tai100\_10\_1 與 tai100\_20\_1 等測資上，NEH 初始化可顯著降低 makespan，顯示其能提供較高品質的初始解。

然而，在少數案例中（如 tai50\_10\_1），NEH 初始化未能帶來明顯優勢，甚至略微劣於隨機初始化。此現象可能與演算法在演化過程中的探索能力有關，當族群已具備足夠多樣性時，初始解品質對最終結果的影響可能較為有限。

整體而言，NEH 初始化能在多數情況下提供較佳的起始解，並加速演算法收斂至高品質解，因此本研究後續實驗將採用 NEH 初始化作為預設設定。
\begin{table}[htbp]
\caption{比較有無NEH初始化之平均與最小Makespan}
\label{tab:neh_min_avg}
\centering
\resizebox{\columnwidth}{!}{
\begin{tabular}{l|cc|cc}
\hline
\textbf{Instance} & \multicolumn{2}{c|}{\textbf{Without NEH}} & \multicolumn{2}{c}{\textbf{With NEH}} \\
 & Min & Avg & Min & Avg \\
\hline
tai20\_5\_1   & 1278 & 1291.30 & 1278 & \textbf{1279.20} \\
tai20\_10\_1  & 1592 & 1607.20 & 1586 & \textbf{1598.00} \\
tai20\_20\_1  & 2305 & 2388.00 & 2312 & \textbf{2327.05} \\
\hline
tai50\_5\_1   & 2724 & 2735.85 & 2724 & \textbf{2728.75} \\
tai50\_10\_1  & 3063 & \textbf{3110.85} & 3072 & 3117.20 \\
tai50\_20\_1  & 3987 & 4019.40 & 3950 & \textbf{3972.10} \\
\hline
tai100\_5\_1  & 5493 & 5499.85 & 5493 & \textbf{5495.75} \\
tai100\_10\_1 & 5850 & 5879.85 & 5805 & \textbf{5825.15} \\
tai100\_20\_1 & 6485 & 6574.80 & 6455 & \textbf{6505.00} \\
\hline
\end{tabular}
}
\end{table}


\subsection{With or without Tabu Mechanism in GA}
% 比較有無 Tabu 機制的 GA (OX) 效能差異

然而，實驗結果顯示，此機制在中大型測資上反而導致解品質顯著下降。進一步分析發現，其原因在於過度強調族群多樣性（over-diversification），削弱了基因演算法在後期進行局部開發（exploitation）的能力。

在原始 GA 中，族群會逐漸收斂至若干高品質區域，並透過交配與微小突變進行細緻調整。然而，當嚴格的 tabu 機制禁止重複或相似個體出現時，演算法會頻繁觸發重新生成或隨機替代個體，導致大量低品質解被引入族群中。

此現象在大型測資中尤其嚴重，因為解空間呈現組合爆炸，使隨機個體與優良解之間距離極大，進而浪費評估資源並破壞原有收斂結構。
\subsection{Effect of Population-level Tabu Mechanism (Static)}

為了分析族群層級禁忌機制（population-level tabu mechanism）對演算法效能之影響，本實驗在固定其他參數設定下，比較是否於基因演算法中啟用靜態 Tabu 機制。具體而言，本實驗固定使用 LOX 作為交配算子、插入突變作為突變方式、突變率設為 0.5，並採用 NEH 初始化；區域搜尋則暫不啟用。比較對象分別為未使用 Tabu 機制之版本與使用固定禁忌長度之靜態 Tabu 版本。

\begin{table}[htbp]
\caption{比較有無static tabu mechanism之平均與最小Makespan}
\label{tab:static_tabu_min_avg}
\centering
\resizebox{\columnwidth}{!}{
\begin{tabular}{l|cc|cc}
\hline
\textbf{Instance} & \multicolumn{2}{c|}{\textbf{Without Tabu}} & \multicolumn{2}{c}{\textbf{With Static Tabu}} \\
 & Min & Avg & Min & Avg \\
\hline
tai20\_5\_1   & 1278 & 1280.20 & 1278 & \textbf{1278.60} \\
tai20\_10\_1  & 1583 & 1600.55 & 1583 & \textbf{1590.45} \\
tai20\_20\_1  & 2316 & 2324.85 & 2306 & \textbf{2322.35} \\
\hline
tai50\_5\_1   & 2729 & 2729.00 & 2729 & 2729.00 \\
tai50\_10\_1  & 3089 & 3121.20 & 3051 & \textbf{3113.30} \\
tai50\_20\_1  & 3950 & 3978.30 & 3962 & \textbf{3977.95} \\
\hline
tai100\_5\_1  & 5493 & 5496.90 & 5493 & \textbf{5496.40} \\
tai100\_10\_1 & 5793 & \textbf{5819.85} & 5799 & 5824.70 \\
tai100\_20\_1 & 6441 & \textbf{6501.75} & 6465 & 6525.55 \\
\hline
\end{tabular}
}
\end{table}

由表~\ref{tab:static_tabu_min_avg} 可觀察到，靜態 Tabu 機制對整體解品質並未帶來一致性的提升。在部分小型測資上，啟用 Tabu 後確實能略微改善平均完工時間，例如 \texttt{tai20\_5\_1} 與 \texttt{tai20\_10\_1} 的平均值皆有所下降；而在 \texttt{tai20\_20\_1} 與 \texttt{tai50\_10\_1} 等案例中，最佳解（Min）也有改善，顯示 Tabu 機制在某些情況下確實能幫助演算法跳脫既有搜尋區域。

然而，這種改善並不穩定。對於部分中大型測資而言，啟用靜態 Tabu 後的平均解品質反而略有退步，例如 \texttt{tai100\_10\_1} 的平均完工時間由 5819.85 上升至 5824.70，\texttt{tai100\_20\_1} 則由 6501.75 上升至 6525.55。此外，在 \texttt{tai50\_10\_1} 中，雖然最佳解由 3089 改善至 3051，但標準差由 13.43 上升至 25.23，代表搜尋結果的穩定性反而下降。這說明靜態 Tabu 機制雖然可能增強探索能力，但未必能穩定地轉化為更好的整體解品質。

從平均代數（AvgGen）與平均執行時間（AvgTime）來看，啟用 Tabu 機制後，大多數測資的世代數與執行時間皆略為增加。例如在 \texttt{tai100\_20\_1} 中，平均代數由 897.00 上升至 912.90，平均執行時間由 3.63 秒增加至 3.91 秒；在 \texttt{tai50\_20\_1} 中，平均執行時間亦由 2.19 秒上升至 2.56 秒。此結果顯示，Tabu 機制雖增加了搜尋過程中的多樣性控制，但同時也提高了計算成本。

此現象可解釋為：靜態 Tabu 機制對重複或相似個體施加了固定強度的限制，雖有助於避免族群過早集中於單一區域，但也可能妨礙演化後期對優良子結構的細緻重組與局部微調。特別是在中大型測資中，當族群逐漸接近高品質解區域時，過度嚴格的禁忌限制可能迫使演算法排除原本具有潛力的相似後代，進而降低 exploitation 能力，使平均結果不升反降。

綜合而言，靜態 Tabu 機制在本研究設定下僅帶來有限且不穩定的改善。雖然其在部分測資上能提升最佳解或平均解品質，但整體而言仍未能穩定優於未使用 Tabu 的版本，且伴隨額外的時間成本。因此，僅以固定禁忌長度控制族群多樣性並不足以有效提升演算法效能，亦顯示後續引入動態調整機制（Dynamic Tabu Tenure）具有必要性。

\clearpage

\subsection{Dynamic Tabu Tenure Mechanism in GA}
% 探討動態禁忌長度 (Dynamic Tabu Tenure) 對效能的恢復與提升
為了解決靜態 Tabu 機制造成的過度多樣化，以及後期局部開發能力不足的問題，本研究進一步實作動態禁忌期限（Dynamic Tabu Tenure）機制。

\textbf{以 FFE 為基礎的動態衰減策略：} \\
本研究觀察到，演化初期需要較高的多樣性以擴大搜尋範圍（exploration），而搜尋後期則需要降低限制，讓族群能對已發現的優良區域進行更細緻的局部開發（exploitation）。因此，本研究不再以世代數作為調整依據，而是改採 FFE 的使用進度來動態控制 tabu tenure。

具體而言，演算法會先計算目前已使用 FFE 佔總 FFE 預算的比例，並依此比例線性遞減禁忌期限。於搜尋初期，由於 FFE 使用比例低，系統維持較長的禁忌期限（初始 $\texttt{tabu\_tenure}=10$），以較嚴格的方式限制重複個體的產生，促進族群探索不同區域；隨著 FFE 持續消耗，禁忌期限逐步縮短；當搜尋接近預算上限時，tabu tenure 下降至接近 0，使演算法能減少多樣性約束，專注於對高品質解進行後期微調。

此設計的優點在於，相較於以世代數作為控制基準，FFE 更能直接反映實際搜尋成本。特別是在不同機制下，每一世代所消耗的評估次數可能並不相同，因此使用 FFE 作為動態調整依據，能更精確地對應演算法當前所處的搜尋階段。

我們進一步比較了三種設定：
未使用 Tabu 機制、
使用固定禁忌長度之靜態 Tabu，
以及使用基於 FFE 動態調整禁忌期限之 Dynamic Tabu。

\begin{table}[htbp]
\caption{Comparison of average makespan and average runtime under different tabu mechanisms}
\label{tab:tabu_avg_time_comp}
\centering
\resizebox{\columnwidth}{!}{
\begin{tabular}{l|cc|cc|cc}
\hline
\textbf{Instance}
& \multicolumn{2}{c|}{\textbf{Without Tabu}}
& \multicolumn{2}{c|}{\textbf{Static Tabu}}
& \multicolumn{2}{c}{\textbf{Dynamic Tabu}} \\
& Avg & AvgTime & Avg & AvgTime & Avg & AvgTime \\
\hline
tai20\_5\_1   & 1279.60 & 0.62 & \textbf{1278.40} & 0.78 & 1279.00 & 0.77 \\
tai20\_10\_1  & 1596.65 & 0.79 & \textbf{1590.75} & 1.00 & 1593.35 & 0.97 \\
tai20\_20\_1  & 2329.40 & 1.09 & 2320.25 & 1.25 & \textbf{2317.90} & 1.21 \\
\hline
tai50\_5\_1   & \textbf{2728.75} & 1.00 & 2729.00 & 1.20 & 2729.00 & 1.18 \\
tai50\_10\_1  & 3117.50 & 1.23 & 3116.30 & 1.43 & \textbf{3108.25} & 1.43 \\
tai50\_20\_1  & \textbf{3970.80} & 2.06 & 3976.75 & 2.35 & 3980.40 & 2.29 \\
\hline
tai100\_5\_1  & 5494.70 & 1.37 & \textbf{5494.20} & 1.58 & 5494.60 & 1.56 \\
tai100\_10\_1 & 5824.35 & 2.04 & \textbf{5823.60} & 2.29 & 5824.85 & 2.29 \\
tai100\_20\_1 & 6512.60 & 3.10 & 6512.55 & 3.41 & \textbf{6510.40} & 3.43 \\
\hline
\end{tabular}
}
\end{table}

實驗結果如表~\ref{tab:tabu_avg_time_comp} 所示，我們進一步比較了未使用 Tabu、使用固定禁忌長度（Static Tabu），以及使用基於 FFE 動態調整禁忌期限（Dynamic Tabu）三種設定在平均解品質與平均執行時間上的差異。

就平均完工時間（Avg）而言，Dynamic Tabu 在部分測資上相較於 Static Tabu 確實展現較佳表現。例如在 \texttt{tai20\_20\_1} 中，平均值由 Static Tabu 的 2320.25 進一步下降至 2317.90；在 \texttt{tai50\_10\_1} 中，亦由 3116.30 改善至 3108.25；而在 \texttt{tai100\_20\_1} 中，Dynamic Tabu 的平均值 6510.40 亦略優於 Static Tabu 的 6512.55。此結果顯示，根據 FFE 使用進度逐步放寬禁忌限制，有助於修正靜態 Tabu 在後期過度抑制 exploitation 的問題。

然而，Dynamic Tabu 並未在所有測資上全面優於其他方法。在 \texttt{tai20\_5\_1} 與 \texttt{tai20\_10\_1} 等小型測資中，Static Tabu 仍具有最佳平均表現；而在 \texttt{tai50\_20\_1} 中，未使用 Tabu 的基準版本反而取得三者中最佳平均值 3970.80。這表示動態禁忌期限雖能改善部分案例中的搜尋平衡，但其效益仍與問題規模及測資特性有關，並非在所有情況下皆優於基準方法。

從平均執行時間（AvgTime）觀察，啟用 Tabu 機制後，不論為 Static 或 Dynamic 版本，其計算時間皆普遍高於未使用 Tabu 的版本。例如在 \texttt{tai100\_20\_1} 中，未使用 Tabu 的平均執行時間為 3.10 秒，而 Static 與 Dynamic 分別增加至 3.41 秒與 3.43 秒；在 \texttt{tai50\_20\_1} 中，亦由 2.06 秒增加至 2.35 秒與 2.29 秒。此結果顯示，Tabu 機制雖然有助於控制族群多樣性並改善部分測資的解品質，但同時也伴隨額外的計算成本。

綜合而言，Dynamic Tabu 相較於 Static Tabu 在若干測資上確實獲得了較佳的平均解品質，顯示以 FFE 作為動態調整依據具有一定效果；然而，從整體表現來看，其優勢尚不足以全面超越未使用 Tabu 的基準版本。換言之，Dynamic Tabu 可視為對 Static Tabu 的有效修正，但其實際效益仍需視問題特性與計算成本一併評估。




\subsection{Effect of Local Search Design}
在確定交配算子、突變設定與 NEH 初始化後，本研究進一步分析 local search 的設計方式對演算法效能之影響。由於 local search 同時涉及作用對象（who）、觸發頻率（when）以及搜尋強度（iters）等多個設計因素，若直接全面搜尋所有組合將造成過高的實驗成本。因此，本研究採分階段方式進行分析：首先比較 local search 的作用對象，其次比較觸發頻率，最後比較搜尋強度。

\subsubsection{Effect of Selection Strategy (Who)}
首先比較 local search 的作用對象，包括僅對當代最佳個體進行 local search（Elite）以及從前五名個體中隨機選取一個進行 local search（Top-5 Random）。如表~\ref{tab:ls_who_short} 所示，Top-5 Random 在多數測資上表現較為穩定，且在部分案例中優於 Elite。以 \texttt{tai50\_10\_1} 為例，Top-5 Random 的平均 makespan 為 3102.70，優於 Elite 的 3107.65；在 \texttt{tai100\_20\_1} 中，Top-5 Random 亦由 6514.30 改善至 6490.50。雖然在某些測資中 Elite 與 Top-5 Random 的差異不大，但整體而言，僅對最佳個體反覆執行 local search 容易導致過度 exploitation，使搜尋較快集中於單一區域；相較之下，從前五名中隨機選取一個進行局部搜尋，能在保留高品質候選解的同時，維持較佳的探索彈性。因此，後續實驗採用 Top-5 Random 作為 local search 的作用對象。


\begin{table}[htbp]
\caption{Local Search選擇不同目標時之平均與最小 Makespan}
\label{tab:ls_who_short}
\centering
\resizebox{\columnwidth}{!}{
\begin{tabular}{l|cc|cc|cc}
\hline
\textbf{Instance}
& \multicolumn{2}{c|}{\textbf{No LS}}
& \multicolumn{2}{c|}{\textbf{Elite LS}}
& \multicolumn{2}{c}{\textbf{Top-5Random LS}} \\
& Min & Avg & Min & Avg & Min & Avg \\
\hline
tai20\_5\_1   & 1278 & 1280.20 & 1278 & 1279.20 & 1278 & 1278.80 \\
tai20\_10\_1  & 1583 & 1595.20 & 1583 & 1599.20 & 1586 & 1595.10 \\
tai20\_20\_1  & 2315 & 2326.05 & 2312 & 2327.95 & 2307 & 2325.20 \\
tai50\_5\_1   & 2729 & 2729.00 & 2729 & 2729.00 & 2729 & 2729.00 \\
tai50\_10\_1  & 3076 & 3114.70 & 3069 & 3107.65 & 3066 & 3102.70 \\
tai50\_20\_1  & 3945 & 3974.80 & 3936 & 3970.80 & 3935 & 3978.10 \\
tai100\_5\_1  & 5493 & 5495.20 & 5493 & 5495.75 & 5493 & 5495.50 \\
tai100\_10\_1 & 5804 & 5823.65 & 5786 & 5820.05 & 5804 & 5821.00 \\
tai100\_20\_1 & 6435 & 6507.90 & 6472 & 6514.30 & 6436 & 6490.50 \\
\hline
\end{tabular}
}
\end{table}


\subsubsection{Effect of Application Frequency (When)}
在固定採用 Top-5 Random 作為作用對象後，本研究進一步比較 local search 的觸發頻率：每 10 代執行一次與每 20 代執行一次。由表~\ref{tab:ls_when_short} 可觀察到，兩者在整體表現上差異有限，而提高 local search 頻率並未帶來穩定改善。例如在 \texttt{tai20\_10\_1} 中，每 20 代執行一次的平均 makespan 為 1595.40，而每 10 代執行一次則上升至 1602.20；在 \texttt{tai50\_10\_1} 中，every 20 的平均值為 3102.00，而 every 10 為 3104.00。雖然在少數大型測資（如 \texttt{tai100\_20\_1}）中 every 10 略優於 every 20，但整體趨勢並不明顯。此結果顯示，local search 並非越頻繁越有效，過於頻繁的局部搜尋會消耗較多 FFE 預算，壓縮 GA 的全域探索能力。因此，本研究後續採用 every 20 generations 作為較穩定的設定。


\begin{table}[htbp]
\caption{不同頻率做Local Search之平均與最小 Makespan}
\label{tab:ls_when_short}
\centering
\resizebox{\columnwidth}{!}{
\begin{tabular}{l|cc|cc}
\hline
\textbf{Instance}
& \multicolumn{2}{c|}{\textbf{Every 10 Generations}}
& \multicolumn{2}{c}{\textbf{Every 20 Generations}} \\
& Min & Avg & Min & Avg \\
\hline
tai20\_5\_1   & 1278 & 1280.40 & 1278 & 1279.20 \\
tai20\_10\_1  & 1583 & 1602.20 & 1582 & 1595.40 \\
tai20\_20\_1  & 2316 & 2329.30 & 2310 & 2329.05 \\
tai50\_5\_1   & 2724 & 2728.75 & 2729 & 2729.10 \\
tai50\_10\_1  & 3056 & 3104.00 & 3050 & 3102.00 \\
tai50\_20\_1  & 3943 & 3972.80 & 3941 & 3972.10 \\
tai100\_5\_1  & 5493 & 5495.85 & 5493 & 5497.50 \\
tai100\_10\_1 & 5793 & 5818.15 & 5785 & 5818.30 \\
tai100\_20\_1 & 6469 & 6502.10 & 6470 & 6508.05 \\
\hline
\end{tabular}
}
\end{table}



\subsubsection{Effect of Search Intensity (iters)}
在固定作用對象為 Top-5 Random 且觸發頻率為 every 20 generations 後，本研究最後比較不同搜尋強度（iters = 20、50、100）的影響。由表~\ref{tab:ls_iters_short} 可知，增加 local search 強度並未帶來一致性的效能提升。在部分案例中，中等強度或較高強度可帶來改善，例如 \texttt{tai100\_20\_1} 中 iters = 50 的平均 makespan 為 6499.15，優於 iters = 20 的 6506.45 與 iters = 100 的 6503.90；然而在 \texttt{tai50\_10\_1} 中，iters = 20 反而最佳，而在 \texttt{tai100\_10\_1} 中，iters = 20 與 50 幾乎相同，iters = 100 則略差。另一方面，由附錄表~\ref{tab:ls_iters_full} 可觀察到，隨著 iters 增加，平均世代數（AvgGen）普遍下降，表示 local search 會消耗更多 FFE 預算，進而壓縮演化世代數。因此，過高的搜尋強度未必能有效提升整體表現，反而可能削弱 GA 的搜尋深度。



\begin{table}[htbp]
\caption{比較不同迭代次數之平均與最小 Makespan}
\label{tab:ls_iters_short}
\centering
\resizebox{\columnwidth}{!}{
\begin{tabular}{l|cc|cc|cc}
\hline
\textbf{Instance}
& \multicolumn{2}{c|}{\textbf{iters = 20}}
& \multicolumn{2}{c|}{\textbf{iters = 50}}
& \multicolumn{2}{c}{\textbf{iters = 100}} \\
& Min & Avg & Min & Avg & Min & Avg \\
\hline
tai20\_5\_1   & 1278 & 1280.00 & 1278 & 1279.60 & 1278 & 1279.60 \\
tai20\_10\_1  & 1583 & 1598.45 & 1586 & 1597.10 & 1586 & 1601.80 \\
tai20\_20\_1  & 2316 & 2329.95 & 2307 & 2324.40 & 2312 & 2327.05 \\
tai50\_5\_1   & 2729 & 2729.00 & 2729 & 2729.00 & 2729 & 2729.00 \\
tai50\_10\_1  & 3051 & 3100.65 & 3056 & 3109.70 & 3054 & 3109.20 \\
tai50\_20\_1  & 3951 & 3974.40 & 3947 & 3973.00 & 3941 & 3975.35 \\
tai100\_5\_1  & 5493 & 5496.90 & 5495 & 5496.55 & 5493 & 5497.75 \\
tai100\_10\_1 & 5800 & 5821.20 & 5800 & 5821.75 & 5797 & 5824.15 \\
tai100\_20\_1 & 6476 & 6506.45 & 6467 & 6499.15 & 6453 & 6503.90 \\
\hline
\end{tabular}
}
\end{table}


\subsubsection{Overall Discussion}

綜合以上三個面向的分析可知，local search 的設計需要在 exploration 與 exploitation 之間取得平衡。若 local search 過度集中於菁英個體、過於頻繁地介入，或設定過高的搜尋強度，皆可能壓縮 GA 的全域探索能力，進而抵銷 local search 所帶來的局部改良效果。在本研究的設定下，較佳的 local search 設計為：從前五名個體中隨機選取一個作為作用對象、每 20 代執行一次，並採用中等強度的搜尋步數（如 iters = 50）。然而整體而言，local search 的改善幅度仍然有限，顯示其效益高度依賴於問題規模與參數設定，並非在所有情況下皆能穩定提升演算法效能。


\clearpage

\section{Conclusion}


\appendices


\section{Detailed Experimental Results}


\begin{table*}[t]
\centering
\caption{Detailed crossover comparison in terms of minimum, average, maximum, and standard deviation}
\label{tab:crossover_full}
\resizebox{\textwidth}{!}{
\begin{tabular}{c|cccc|cccc|cccc|cccc}
\hline
Instance 
& \multicolumn{4}{c|}{OX} 
& \multicolumn{4}{c|}{LOX} 
& \multicolumn{4}{c|}{PMX} 
& \multicolumn{4}{c}{CX} \\
\hline
& Min & Avg & Max & Std
& Min & Avg & Max & Std
& Min & Avg & Max & Std
& Min & Avg & Max & Std \\
\hline
tai20\_5\_1   & 1278 & 1295.85 & 1302 & 4.82 & 1278 & 1296.05 & 1297 & 4.25 & 1279 & 1294.60 & 1302 & 6.68 & 1278 & 1297.10 & 1305 & 5.20 \\
tai20\_10\_1  & 1608 & 1627.55 & 1651 & 11.03 & 1605 & 1621.85 & 1654 & 12.76 & 1599 & 1623.40 & 1653 & 13.67 & 1593 & 1629.55 & 1671 & 22.13 \\
tai20\_20\_1  & 2331 & 2373.00 & 2430 & 32.31 & 2318 & 2376.05 & 2446 & 35.80 & 2317 & 2360.80 & 2435 & 33.68 & 2323 & 2363.25 & 2417 & 30.51 \\
tai50\_5\_1   & 2724 & 2742.85 & 2774 & 12.39 & 2724 & 2741.35 & 2774 & 11.51 & 2724 & 2742.80 & 2774 & 12.09 & 2729 & 2744.75 & 2774 & 11.42 \\
tai50\_10\_1  & 3098 & 3145.70 & 3206 & 30.63 & 3114 & 3152.65 & 3246 & 31.22 & 3109 & 3150.80 & 3213 & 25.52 & 3102 & 3150.05 & 3207 & 27.49 \\
tai50\_20\_1  & 4016 & 4056.80 & 4100 & 23.22 & 4013 & 4052.85 & 4121 & 30.92 & 4002 & 4054.40 & 4111 & 33.06 & 3993 & 4057.50 & 4132 & 33.81 \\
tai100\_5\_1  & 5495 & 5512.90 & 5561 & 17.02 & 5493 & 5517.20 & 5564 & 19.81 & 5495 & 5516.15 & 5564 & 17.44 & 5493 & 5513.15 & 5536 & 14.51 \\
tai100\_10\_1 & 5888 & 5963.00 & 6057 & 47.37 & 5877 & 5939.00 & 6018 & 38.70 & 5884 & 5970.60 & 6051 & 46.37 & 5868 & 5959.85 & 6030 & 38.62 \\
tai100\_20\_1 & 6584 & 6687.35 & 6776 & 45.03 & 6589 & 6685.45 & 6807 & 54.39 & 6582 & 6674.90 & 6771 & 52.06 & 6598 & 6679.70 & 6764 & 44.14 \\
\hline
\end{tabular}
}
\end{table*}



\begin{table*}[t]
\centering
\caption{Detailed mutation setting comparison in terms of minimum, average, maximum, and standard deviation}
\label{tab:mutation_settings_full}
\resizebox{\textwidth}{!}{
\begin{tabular}{c|cccc|cccc|cccc|cccc}
\hline
Instance 
& \multicolumn{4}{c|}{Insertion+0.1}
& \multicolumn{4}{c|}{Insertion+0.5}
& \multicolumn{4}{c|}{Swap+0.1}
& \multicolumn{4}{c}{Swap+0.5} \\
\hline
& Min & Avg & Max & Std
& Min & Avg & Max & Std
& Min & Avg & Max & Std
& Min & Avg & Max & Std \\
\hline
tai20\_5\_1   & 1278 & 1293.35 & 1297 & 7.51 & 1278 & 1291.30 & 1297 & 8.93 & 1297 & 1297.80 & 1305 & 2.46 & 1297 & 1297.00 & 1297 & 0.00 \\
tai20\_10\_1  & 1603 & 1624.35 & 1657 & 13.96 & 1583 & 1603.65 & 1621 & 12.64 & 1604 & 1640.80 & 1702 & 25.24 & 1586 & 1617.65 & 1661 & 15.88 \\
tai20\_20\_1  & 2312 & 2358.40 & 2424 & 31.29 & 2312 & 2341.05 & 2400 & 23.94 & 2334 & 2381.40 & 2433 & 30.06 & 2307 & 2346.30 & 2401 & 22.79 \\
tai50\_5\_1   & 2731 & 2743.70 & 2752 & 7.13 & 2724 & 2734.80 & 2752 & 7.60 & 2729 & 2743.75 & 2752 & 10.32 & 2724 & 2729.20 & 2752 & 6.42 \\
tai50\_10\_1  & 3117 & 3145.95 & 3198 & 23.10 & 3082 & 3116.25 & 3149 & 18.67 & 3126 & 3157.45 & 3227 & 22.59 & 3107 & 3139.20 & 3170 & 20.60 \\
tai50\_20\_1  & 4021 & 4073.40 & 4173 & 40.80 & 3986 & 4011.25 & 4046 & 17.65 & 4038 & 4089.00 & 4139 & 33.33 & 4001 & 4051.35 & 4127 & 31.85 \\
tai100\_5\_1  & 5493 & 5511.75 & 5546 & 16.76 & 5493 & 5496.10 & 5527 & 7.75 & 5495 & 5512.20 & 5538 & 16.43 & 5495 & 5506.80 & 5529 & 15.59 \\
tai100\_10\_1 & 5898 & 5962.25 & 6062 & 43.20 & 5823 & 5859.90 & 5897 & 23.01 & 5907 & 5993.90 & 6054 & 39.67 & 5850 & 5906.55 & 5959 & 32.91 \\
tai100\_20\_1 & 6563 & 6685.95 & 6753 & 50.59 & 6471 & 6579.90 & 6686 & 49.40 & 6580 & 6690.70 & 6811 & 59.00 & 6500 & 6581.90 & 6713 & 46.24 \\
\hline
\end{tabular}
}
\end{table*}



\begin{table*}[t]
\caption{Detailed comparison with and without NEH initialization}
\label{tab:neh_full}
\centering
\resizebox{\textwidth}{!}{
\begin{tabular}{l|cccc|cccc}
\hline
\textbf{Instance} & \multicolumn{4}{c|}{\textbf{Without NEH}} & \multicolumn{4}{c}{\textbf{With NEH}} \\
 & Min & Avg & Max & Std & Min & Avg & Max & Std \\
\hline
tai20\_5\_1   & 1278 & 1291.30 & 1297 & 8.93  & 1278 & 1279.20 & 1286 & 2.93  \\
tai20\_10\_1  & 1592 & 1607.20 & 1625 & 11.24 & 1586 & 1598.00 & 1618 & 11.17 \\
tai20\_20\_1  & 2305 & 2338.00 & 2388 & 23.71 & 2312 & 2327.05 & 2345 & 10.43 \\
\hline
tai50\_5\_1   & 2724 & 2735.85 & 2752 & 6.43  & 2724 & 2728.75 & 2729 & 1.12  \\
tai50\_10\_1  & 3063 & 3110.85 & 3168 & 27.35 & 3072 & 3117.20 & 3138 & 22.01 \\
tai50\_20\_1  & 3987 & 4019.40 & 4050 & 17.78 & 3950 & 3972.10 & 4000 & 15.48 \\
\hline
tai100\_5\_1  & 5493 & 5499.85 & 5527 & 11.60 & 5493 & 5495.75 & 5514 & 5.53  \\
tai100\_10\_1 & 5850 & 5879.85 & 5934 & 24.01 & 5805 & 5825.15 & 5869 & 16.58 \\
tai100\_20\_1 & 6485 & 6574.80 & 6632 & 43.70 & 6455 & 6505.00 & 6564 & 30.06 \\
\hline
\end{tabular}
}
\end{table*}



\begin{table*}[t]
\caption{Detailed comparison with and without static tabu mechanism}
\label{tab:static_tabu_full}
\centering
\resizebox{\textwidth}{!}{
\begin{tabular}{l|cccccc|cccccc}
\hline
\textbf{Instance} & \multicolumn{6}{c|}{\textbf{Without Tabu}} & \multicolumn{6}{c}{\textbf{With Static Tabu}} \\
 & Min & Avg & Max & Std & AvgGen & AvgTime & Min & Avg & Max & Std & AvgGen & AvgTime \\
\hline
tai20\_5\_1   & 1278 & 1280.20 & 1286 & 3.55  & 995.00 & 0.66 & 1278 & 1278.60 & 1286 & 1.96  & 1013.00 & 0.81 \\
tai20\_10\_1  & 1583 & 1600.55 & 1626 & 11.32 & 995.00 & 0.83 & 1583 & 1590.45 & 1601 & 5.74  & 1012.20 & 1.02 \\
tai20\_20\_1  & 2316 & 2324.85 & 2340 & 6.76  & 995.00 & 1.14 & 2306 & 2322.35 & 2335 & 8.52  & 1012.20 & 1.29 \\
tai50\_5\_1   & 2729 & 2729.00 & 2729 & 0.00  & 973.00 & 1.00 & 2729 & 2729.00 & 2729 & 0.00  & 991.30  & 1.19 \\
tai50\_10\_1  & 3089 & 3121.20 & 3138 & 13.43 & 973.00 & 1.27 & 3051 & 3113.30 & 3135 & 25.23 & 990.60  & 1.50 \\
tai50\_20\_1  & 3950 & 3978.30 & 3998 & 14.50 & 973.00 & 2.19 & 3962 & 3977.95 & 3998 & 10.72 & 990.15  & 2.56 \\
tai100\_5\_1  & 5493 & 5496.90 & 5514 & 6.69  & 897.00 & 1.38 & 5493 & 5496.40 & 5514 & 6.27  & 913.65  & 1.58 \\
tai100\_10\_1 & 5793 & 5819.85 & 5837 & 13.77 & 897.00 & 1.92 & 5799 & 5824.70 & 5845 & 14.48 & 913.00  & 2.14 \\
tai100\_20\_1 & 6441 & 6501.75 & 6560 & 36.96 & 897.00 & 3.63 & 6465 & 6525.55 & 6563 & 22.61 & 912.90  & 3.91 \\
\hline
\end{tabular}
}
\end{table*}



\begin{table*}[t]
\caption{Detailed comparison of tabu mechanisms}
\label{tab:tabu_full_comp}
\centering
\resizebox{\textwidth}{!}{
\begin{tabular}{l|cccccc|cccccc|cccccc}
\hline
\textbf{Instance}
& \multicolumn{6}{c|}{\textbf{Without Tabu}}
& \multicolumn{6}{c|}{\textbf{Static Tabu}}
& \multicolumn{6}{c}{\textbf{Dynamic Tabu}} \\
& Min & Avg & Max & Std & AvgGen & AvgTime
& Min & Avg & Max & Std & AvgGen & AvgTime
& Min & Avg & Max & Std & AvgGen & AvgTime \\
\hline
tai20\_5\_1
& 1278 & 1279.60 & 1286 & 3.02 & 995.00 & 0.62
& 1278 & 1278.40 & 1282 & 1.23 & 1013.00 & 0.78
& 1278 & 1279.00 & 1286 & 2.55 & 1013.00 & 0.77 \\

tai20\_10\_1
& 1583 & 1596.65 & 1611 & 8.72 & 995.00 & 0.79
& 1583 & 1590.75 & 1613 & 9.07 & 1012.15 & 1.00
& 1586 & 1593.35 & 1617 & 8.93 & 1012.50 & 0.97 \\

tai20\_20\_1
& 2310 & 2329.40 & 2346 & 9.21 & 995.00 & 1.09
& 2309 & 2320.25 & 2348 & 9.33 & 1012.35 & 1.25
& 2307 & 2317.90 & 2337 & 7.49 & 1012.40 & 1.21 \\

\hline
tai50\_5\_1
& 2724 & 2728.75 & 2729 & 1.12 & 973.00 & 1.00
& 2729 & 2729.00 & 2729 & 0.00 & 991.25 & 1.20
& 2729 & 2729.00 & 2729 & 0.00 & 991.45 & 1.18 \\

tai50\_10\_1
& 3066 & 3117.50 & 3138 & 20.94 & 973.00 & 1.23
& 3035 & 3116.30 & 3131 & 23.78 & 990.50 & 1.43
& 3051 & 3108.25 & 3131 & 28.55 & 990.70 & 1.43 \\

tai50\_20\_1
& 3946 & 3970.80 & 3998 & 14.31 & 973.00 & 2.06
& 3947 & 3976.75 & 3993 & 12.67 & 990.35 & 2.35
& 3952 & 3980.40 & 4004 & 15.36 & 990.40 & 2.29 \\

\hline
tai100\_5\_1
& 5493 & 5494.70 & 5505 & 2.79 & 897.00 & 1.37
& 5493 & 5494.20 & 5505 & 2.71 & 913.80 & 1.58
& 5493 & 5494.60 & 5505 & 2.64 & 913.80 & 1.56 \\

tai100\_10\_1
& 5794 & 5824.35 & 5846 & 15.83 & 897.00 & 2.04
& 5795 & 5823.60 & 5837 & 14.73 & 913.00 & 2.29
& 5800 & 5824.85 & 5845 & 15.10 & 913.00 & 2.29 \\

tai100\_20\_1
& 6463 & 6512.60 & 6550 & 23.56 & 897.00 & 3.10
& 6475 & 6512.55 & 6547 & 21.58 & 912.90 & 3.41
& 6473 & 6510.40 & 6543 & 19.20 & 912.75 & 3.43 \\

\hline
\end{tabular}
}
\end{table*}



\end{document}
